\section{Conclusion}

The datum and viscosity determine one pressure-normalized velocity history and
one compatible family of finite pressure-complete rectangles.  Their
mixed-index exhaustion gives every temporal--spatial Sobolev row.  The
completed critical-height identities are evaluated on those same rectangles;
their height exhaustion gives every spatial Sobolev row, and the equation
reconstructs the same temporal and pressure rows.  The scalar
\(\mathfrak R_{N,M}\) measures \(\cR_{N,M}\),
Proposition~\ref{u:prop:datum-terminal-realization} identifies its terminal
value, and Theorem~\ref{u:thm:compatible-inverse-intersection} proves
that every larger rectangle restricts to the same rows in each smaller one and
that the family exhausts all finite mixed indices.

The adjacent rows in each finite rectangle give continuous Sobolev traces
through the proposed terminal time.  Their compatible endpoint values determine
\(u_*\in H^\infty_\sigma\), and the pressure formula and velocity
recurrence recover its complete jet.  Local well-posedness restarts from
that state, while uniqueness on the overlap joins the continuation to the
incoming history.  A finite maximal time would admit an extension, so the
maximal interval is \([0,\infty)\).

These same rectangles also quantify the critical-height rows, frequency tails,
Lipschitz action, and relative-energy stability, with constants depending on
the datum, viscosity, and horizon.  In the small-critical regime, the
scale-invariant norm closes the critical stock/action estimate, and the higher
rectangles admit explicit formulas using the corresponding finite collection
of higher datum norms.

The resulting global solution is smooth, satisfies the exact energy identity,
and carries the decay-normalized pressure.  Relative energy then identifies
every Leray--Hopf solution with that datum with this smooth velocity, and
the pressure is unique up to the unavoidable time-dependent gauge unless this decay normalization is imposed.  The endpoint contradiction and these
consequences all belong to the one fixed-viscosity evolution constructed from
\(u_0\).
